\documentclass{tietreport}

\usepackage{graphicx}
\usepackage{amsmath}

% Bibliography
\usepackage[
  style=alphabetic,
  backend=biber,
  sorting=nty,
  doi=false,
  isbn=false,
  eprint=false
]{biblatex}
\addbibresource{references.bib}

% Chapters start on the next page (no blank pages)
\AtBeginDocument{\let\cleardoublepage\clearpage}

%% ================= PROJECT DETAILS =================

\TitlePageHeader{UCS503 - Software Engineering (Project)}
\ProjectTitle{SafeRoute}
\ProjectSubTitle{A Crowdsourced Safety-Aware Navigation Web App}

\author{%
  \texttt{1024030176} Dhriti Handa \\
  \texttt{1024030180} Disha Kundra%
}

\TitlePageSubText{Group: \texttt{3C16} BE Third Year, COE}
\AdvisorName{Dr. Paramveer Sidhu}

\TitlePageFooterText{{%
    \large Mid-Semester Evaluation \\[0.2cm]
    \bfseries Computer Science and Engineering Department%
  } \\[0.15cm]
  Thapar Institute of Engineering and Technology, Patiala%
}

%% ===================================================

\begin{document}

\maketitle
\tableofcontents

% =====================================================
\chapter{Introduction}
% =====================================================

\section{Background and Context}

When we walk home at night, we do not only ask ``which road is
shortest?'' but also ``which road feels safe?''. Street lighting,
foot traffic and past incidents can differ a lot between two roads
that take the same time to walk. Today this choice is made by
word of mouth, not by any tool.

\textbf{SafeRoute} is an install-free web app that shows two routes
between a start and a destination: the \emph{fastest} route and the
\emph{safest} route. Safety comes from crowdsourced reports, and the
safety score is higher at night.

\textbf{Guiding principle (time-to-value).} Our proposal says to show
a working end-to-end loop as early as possible. So the prototype uses
manually seeded data and a light technology stack, and we add the
heavier parts later.

\subsection{Problem Statement}

Google Maps and Apple Maps optimise only for time, distance and
traffic. They give no safety information. The problem has three parts:

\begin{itemize}
  \item \textbf{No safety option:} the shortest path may go through a
    dark or isolated street and the user cannot know in advance.
  \item \textbf{Safety and SOS apps do not navigate:} Safetipin
    collects safety data but is a planning tool for city governments;
    bSafe and Noonlight send SOS alerts but give no route advice.
  \item \textbf{No time-of-day awareness:} a street that is fine at
    6~pm can feel very different at 11~pm.
\end{itemize}

\section{Existing Solutions and the Gap}

Table~\ref{tab:existing} compares existing tools with SafeRoute.
SafeRoute is the only one that combines walking routes, safety data
and SOS in one install-free web app. (The SafeRoute row is the
intended product; Chapter~3 states what the prototype already does.)

\begin{table}[h]
\centering
\small
\begin{tabular}{|l|c|c|c|c|c|}
  \hline
  \textbf{Tool} & \textbf{Walking} & \textbf{Safety} & \textbf{SOS} &
  \textbf{No} & \textbf{Time of} \\
  & \textbf{routing} & \textbf{data} & & \textbf{install} & \textbf{day} \\
  \hline
  Google / Apple Maps & Yes & No  & No  & Yes & No  \\ \hline
  Safetipin           & No  & Yes & No  & No  & No  \\ \hline
  bSafe / Noonlight   & No  & No  & Yes & No  & No  \\ \hline
  SafeRoute           & Yes & Yes & Yes & Yes & Yes \\ \hline
\end{tabular}
\caption{Comparison with existing tools}
\label{tab:existing}
\end{table}

\section{Project Objectives}

\begin{enumerate}
  \item \textbf{Safety-aware routing:} compute a fastest route
    (distance only) and a safest route (distance weighted by risk).
  \item \textbf{Crowdsourced reporting:} users submit geotagged
    reports with an optional photo; each report raises the risk of the
    nearest road segment.
  \item \textbf{Time-of-day awareness:} risk is increased at night, so
    routes and the heatmap change after dark.
  \item \textbf{Emergency sharing:} an SOS feature creates a
    live-location link for a saved contact.
  \item \textbf{Measurable evaluation:} the safest route should differ
    from the fastest route in at least 80\% of test cases that contain
    a low-safety segment.
\end{enumerate}

\section{Scope and Deliverables}

Table~\ref{tab:status} shows the status of each deliverable at this
mid-semester evaluation.

\begin{table}[h]
\centering
\small
\begin{tabular}{|p{6cm}|l|p{5cm}|}
  \hline
  \textbf{Deliverable} & \textbf{Status} & \textbf{Note} \\ \hline
  Repository and architecture setup & Done & Backend + frontend \\ \hline
  Reporting and storage schema & Done & SQLite (not PostGIS) \\ \hline
  Fastest and safest routing & Done & Custom Dijkstra, seeded data \\ \hline
  Heatmap and night mode & Done & Backend endpoint \\ \hline
  SOS quick-share & Partial & SMS is simulated \\ \hline
  Web UI (home, report, SOS) & Done & Placeholder data \\ \hline
  Frontend connected to API & Pending & Next task \\ \hline
  Real map (Mapbox), IDW interpolation & Pending & Planned \\ \hline
  Push alerts, CI/CD, load test, deployment & Pending & Later milestones \\ \hline
\end{tabular}
\caption{Deliverable status at mid-semester}
\label{tab:status}
\end{table}

% =====================================================
\chapter{Methodology and Design}
% =====================================================

\section{System Architecture}

\subsection{High-Level Design}

SafeRoute has three layers. A React web client talks to a Flask REST
API over HTTP/JSON. The API calls two Python modules (one for routing,
one for safety) and stores everything in an SQLite database.

\begin{figure}[h]
\centering
\fbox{\parbox{0.8\textwidth}{\centering
  \textbf{React web client (Vite)} \\
  Home, route and heatmap, report, SOS}}\\[3pt]
$\downarrow$ HTTP / JSON \\[3pt]
\fbox{\parbox{0.8\textwidth}{\centering
  \textbf{Flask REST API} (\texttt{app.py})}}\\[3pt]
$\downarrow$ \\[3pt]
\fbox{\parbox{0.38\textwidth}{\centering
  \textbf{routing.py} \\ Dijkstra, two modes}}\hspace{0.02\textwidth}
\fbox{\parbox{0.38\textwidth}{\centering
  \textbf{safety.py} \\ risk, heatmap, reports}}\\[3pt]
$\downarrow$ \\[3pt]
\fbox{\parbox{0.8\textwidth}{\centering
  \textbf{SQLite database} + \texttt{uploads/} for photos}}
\caption{Architecture of the SafeRoute prototype}
\label{fig:arch}
\end{figure}

\subsection{Core Workflow}

\begin{enumerate}
  \item The user picks a start and a destination
    (\texttt{GET /api/places}).
  \item The app asks for routes (\texttt{GET /api/route}); the backend
    returns the fastest and the safest route.
  \item If the user sees a hazard, they send a report
    (\texttt{POST /api/reports}).
  \item The report raises the risk of that road, so the heatmap and
    the next route request show it immediately.
\end{enumerate}

\subsection{Technical Stack}

Python and Flask~\cite{flask} for the backend, React~\cite{react} with
Vite~\cite{vite} for the frontend, and SQLite~\cite{sqlite} for the
database. Table~\ref{tab:stack} compares this with the stack planned in
the proposal. We chose the lighter stack first so that the full loop
could be shown early.

\begin{table}[h]
\centering
\small
\begin{tabular}{|l|l|l|}
  \hline
  \textbf{Layer} & \textbf{Prototype (done)} & \textbf{Proposal (planned)} \\ \hline
  Frontend  & React, Vite & + Mapbox GL JS \\ \hline
  Database  & SQLite & PostgreSQL + PostGIS \\ \hline
  Routing   & Own Dijkstra & OSRM~\cite{luxen2011} \\ \hline
  Gap filling & None & IDW interpolation~\cite{shepard1968} \\ \hline
  SOS       & Simulated SMS & Real SMS / push \\ \hline
  Testing   & Manual checks & pytest, Jest, GitHub Actions \\ \hline
\end{tabular}
\caption{Prototype stack versus planned stack}
\label{tab:stack}
\end{table}

\subsection{Design Decisions}

\begin{itemize}
  \item \textbf{Explainable scoring:} one simple formula decides the
    route, so it is easy to test and to explain.
  \item \textbf{Stateless API:} all state is in the database, so the
    API can be scaled later.
  \item \textbf{Never block the user:} routes are ranked by safety but
    never hidden; the user chooses.
\end{itemize}

\section{Data Design}

The database has five tables (Table~\ref{tab:schema}). A road segment is
an undirected \emph{edge} between two places (\emph{nodes}).

\begin{table}[h]
\centering
\small
\begin{tabular}{|l|p{9cm}|}
  \hline
  \textbf{Table} & \textbf{Main columns} \\ \hline
  \texttt{nodes}    & code, name, lat, lng \\ \hline
  \texttt{edges}    & id, u, v, name, risk (0 to 1), note \\ \hline
  \texttt{reports}  & id, edge\_id, lat, lng, categories, description, photo, created \\ \hline
  \texttt{contacts} & id, name, phone \\ \hline
  \texttt{sos}      & id, contact\_id, lat, lng, active, updated \\ \hline
\end{tabular}
\caption{Database schema}
\label{tab:schema}
\end{table}

\subsection{Seeded Campus Data}

\texttt{seed.py} fills the database with 8 places around the Thapar
campus (G1 = Thapar Gate~1, A = Bus Stand Rd Jn, B = Library Block,
C = Model Town Chowk, D = Station Rd Jn, E = Campus Rd Jn,
F = Campus Rd North, BUS = Rajpura Bus Stand) and the 9 road segments
in Table~\ref{tab:edges}, with hand-assigned risk values. This is the
manual seeding the proposal plans for the cold-start phase.

\begin{table}[h]
\centering
\small
\begin{tabular}{|c|l|l|c|p{4.6cm}|}
  \hline
  \textbf{ID} & \textbf{Segment} & \textbf{Road} & \textbf{Risk} &
  \textbf{Note} \\ \hline
  1 & G1 -- A   & Bus Stand Rd  & 0.10 & Well-lit main road \\ \hline
  2 & A -- C    & Model Town Ln & 0.80 & Dim stretch \\ \hline
  3 & C -- D    & Model Town Ln & 0.70 & Isolated, no foot traffic \\ \hline
  4 & D -- BUS  & Station Rd    & 0.30 & Mixed reports \\ \hline
  5 & A -- B    & Bus Stand Rd  & 0.10 & Well-lit, foot traffic \\ \hline
  6 & B -- E    & Campus Rd     & 0.10 & Verified safe \\ \hline
  7 & E -- F    & Campus Rd     & 0.15 & Verified safe \\ \hline
  8 & F -- BUS  & Campus Rd     & 0.10 & Verified safe, shops open \\ \hline
  9 & B -- D    & Link Ln       & 0.50 & Poorly lit after 8 pm \\ \hline
\end{tabular}
\caption{Seeded road segments}
\label{tab:edges}
\end{table}

\section{Implementation Details}

\subsection{Module 1: Safety-Weighted Routing}

\paragraph{Step 1: night adjustment.} At night the stored risk $r$ is
multiplied by $1.25$ (maximum $1$):
\[
  r_{\mathrm{eff}} =
  \begin{cases}
    \min(1,\ 1.25\,r) & \text{night mode} \\
    r & \text{otherwise}
  \end{cases}
\]
The risk is then labelled: below $0.3$ is \emph{safe}, from $0.3$ to
$0.6$ is \emph{mixed}, and $0.6$ or more is \emph{risk}.

\paragraph{Step 2: cost of a segment.} With $d$ the length of the
segment in metres (haversine formula on the coordinates):
\[
  w =
  \begin{cases}
    d & \text{fastest route} \\
    d\,(1 + 4\,r_{\mathrm{eff}}) & \text{safest route}
  \end{cases}
\]
A fully risky segment costs five times its length, so the safest route
accepts a small detour to avoid it; a segment with no risk costs
exactly its length.

\paragraph{Step 3: shortest path.} Dijkstra's algorithm~\cite{dijkstra1959}
(with a binary heap) runs twice on the road graph, once with each cost.
The API returns, for each route, the coordinates, distance, walking time
(83~m/min), average risk and turn-by-turn steps, plus a
\texttt{diverges} flag that says whether the two routes differ.

\subsection{Module 2: Secondary Features}

\paragraph{REST API.} Table~\ref{tab:api} lists the endpoints.

\begin{table}[h]
\centering
\small
\begin{tabular}{|l|p{4.6cm}|p{5.6cm}|}
  \hline
  \textbf{Method} & \textbf{Endpoint} & \textbf{Purpose} \\ \hline
  GET  & \texttt{/api/places}  & List places \\ \hline
  GET  & \texttt{/api/route}   & Fastest and safest route \\ \hline
  GET  & \texttt{/api/heatmap} & Risk of every segment \\ \hline
  POST & \texttt{/api/reports} & Submit a report (+ photo) \\ \hline
  GET  & \texttt{/api/reports} & Latest 20 reports \\ \hline
  GET/POST & \texttt{/api/contacts} & List or add contacts \\ \hline
  POST & \texttt{/api/sos}     & Start an SOS session \\ \hline
  GET/PUT/DELETE & \texttt{/api/sos/<id>} & Read, update or end session \\ \hline
\end{tabular}
\caption{Backend REST endpoints}
\label{tab:api}
\end{table}

\paragraph{Incident reporting.} A report has a location, categories, a
description and an optional photo (saved with a random filename). The
backend finds the nearest road segment, raises its risk by $0.15$
(maximum $1$) and updates its note. Since routing reads the same table,
the next route request already uses the new risk.

\paragraph{Heatmap.} \texttt{/api/heatmap} returns every segment with
its effective risk; with \texttt{night=1} the same map looks riskier.

\paragraph{SOS quick-share.} Starting SOS creates an 8-character session
ID for a saved contact and stores the user's location. The location can
be updated and the session ended later. In the prototype the SMS is
simulated: the tracking link is printed in the server console.

\paragraph{Frontend.} The React app has three screens: \emph{Home}
(places, fastest/safest toggle, route and heatmap tabs), \emph{Report}
(issue types and details) and \emph{SOS} (hold button, share location,
stand down). Currently they use placeholder data and a drawn map;
connecting them to the API is the main remaining task.

\paragraph{Running the prototype.} \texttt{python app.py} in
\texttt{code/backend} (the database is seeded automatically) and
\texttt{npm install \&\& npm run dev} in \texttt{code/frontend}.

% =====================================================
\chapter{Results and Evaluation}
% =====================================================

\section{Evaluation Criteria}

Our main metric is \textbf{route divergence}: the safest route should
differ from the fastest in at least 80\% of test cases where the fastest
route has a risky segment. Secondary metrics are listed in
Table~\ref{tab:criteria}.

\begin{table}[h]
\centering
\small
\begin{tabular}{|p{5cm}|l|p{5cm}|}
  \hline
  \textbf{Metric} & \textbf{Target} & \textbf{Status} \\ \hline
  Route divergence & $\ge 80\%$ & Measured: 43\% \\ \hline
  Report-to-heatmap time & $< 10$ s & Measured: a few ms (in-process) \\ \hline
  Interpolation check, load test, uptime & -- & Pending \\ \hline
\end{tabular}
\caption{Evaluation criteria and status}
\label{tab:criteria}
\end{table}

\section{Testing Methodology}

\begin{itemize}
  \item \textbf{Unit testing:} routing and safety functions were run
    directly in Python on the seeded data; automated \texttt{pytest}
    tests are planned.
  \item \textbf{Integration testing:} \texttt{/api/places} and
    \texttt{/api/route} were called in a browser and returned the
    expected data. Report, contact and SOS endpoints are not yet tested
    end to end.
  \item \textbf{Divergence testing:} the router was run for all
    $8 \times 7 = 56$ ordered pairs of places in night mode.
  \item \textbf{Report test:} one report was added and routes were
    recomputed.
  \item \textbf{Performance and user testing:} timing was measured
    in-process only; load test and user pilot come later.
\end{itemize}

\section{Performance Metrics}

\subsection{Headline Route}

For Thapar Gate~1 to Rajpura Bus Stand, the fastest route uses Link Lane
(poorly lit) and Station Road (mixed). The safest route stays on Campus
Road. Table~\ref{tab:routes} shows the result.

\begin{table}[h]
\centering
\begin{tabular}{|l|l|r|r|r|}
  \hline
  \textbf{Setting} & \textbf{Route} & \textbf{Distance (m)} &
  \textbf{Time (min)} & \textbf{Avg.\ risk} \\ \hline
  Day   & Fastest & 976  & 12 & 0.26 \\ \hline
  Day   & Safest  & 1030 & 12 & 0.11 \\ \hline
  Night & Fastest & 976  & 12 & 0.32 \\ \hline
  Night & Safest  & 1030 & 12 & 0.14 \\ \hline
\end{tabular}
\caption{Fastest versus safest route, G1 to BUS}
\label{tab:routes}
\end{table}

The safest route is only 54~m (about 5.5\%) longer but cuts the average
risk by more than half (58\% by day, 56\% at night).

\subsection{Route Divergence}

A pair of places is a test case if its fastest route has at least one
\emph{risk} segment. There were 28 such cases, and the routes diverged
in 12 of them, i.e.\ \textbf{43\%}, below the 80\% target.

\subsection{Effect of a Report}

A ``Poor lighting'' report on Link Lane raised its risk from 0.50 to
0.65. The fastest route's night risk rose from 0.32 to 0.36, while the
safest route stayed at 0.14 because it already avoids Link Lane.

\subsection{Query Time}

Computing both routes took about 0.13~ms on average over 500 runs. This
is in-process on a tiny graph (8 places), so it is not a substitute for
the planned load test.

\section{Analysis}

\begin{itemize}
  \item \textbf{Did we meet the objectives?} Safety-aware routing,
    reporting, night mode and SOS work in the backend. The 80\%
    divergence target was \emph{not} met (43\%).
  \item \textbf{Why is divergence low?} Of the 16 cases that did not
    diverge, 8 involve Model Town Chowk (C), which is connected only
    through two risky roads, so no safer path exists. The other 8
    involve Station Rd Jn (D): the safer path exists but is much longer,
    so the cost still prefers the direct road. Raising the risk weight
    from 4 to 8, 12 and 20 did not change the result, so the limit
    comes from the small map, not the formula.
  \item \textbf{What worked well?} The cost function is simple, one
    algorithm serves both modes, and reports and night mode change the
    output immediately.
  \item \textbf{Challenges:} PostGIS and OSRM were too heavy for the
    first iteration, so we built the cost function over a small graph.
    Frontend and backend were built in parallel, so integration is left
    for the next phase.
  \item \textbf{Known gaps:} frontend not connected to the API; map is a
    drawing; a report adds a fixed $+0.15$ with no upvotes; any request
    can raise a risk (no protection against fake reports).
\end{itemize}

\section{Risks and Mitigations}

\begin{table}[h]
\centering
\small
\begin{tabular}{|p{3.2cm}|p{4.6cm}|p{4.6cm}|}
  \hline
  \textbf{Risk} & \textbf{Mitigation} & \textbf{Status} \\ \hline
  No data for new areas (cold start) & Seed scores; interpolate later &
    Seeded; interpolation pending \\ \hline
  Fake reports & Rate limit, minimum upvotes & Not implemented \\ \hline
  OSRM complexity & Own router first & Avoided for now \\ \hline
  Slow queries at scale & Spatial indexes, load test & Pending \\ \hline
\end{tabular}
\caption{Main risks from the proposal}
\label{tab:risks}
\end{table}

% =====================================================
\chapter{Conclusion}
% =====================================================

\section{Summary of Work}

We built a Flask backend with a seeded campus graph, a safety-weighted
Dijkstra router (fastest and safest), night-time risk adjustment, a
heatmap, photo-supported reporting that updates road risk, and a
simulated SOS flow. We also built the React screens. The routing,
reporting, time-of-day and SOS objectives work in the backend; the
divergence target needs a larger map.

\section{Key Findings}

\begin{itemize}
  \item A simple cost $d(1+4r)$ steers routes away from risky roads: a
    54~m detour halved the average risk on the headline route.
  \item A report affects the very next route request, because routing
    and reporting share one table.
  \item Divergence depends on map structure: only 43\% on the small
    seeded graph, since it has few parallel streets.
\end{itemize}

\section{Future Work and Recommendations}

\begin{enumerate}
  \item Connect the React screens to the API and add a real map.
  \item Move to PostgreSQL with PostGIS; compare with OSRM routing.
  \item Add IDW interpolation~\cite{shepard1968} for areas without
    reports, and weight reports by recency and upvotes.
  \item Add login, rate limiting and automated tests (\texttt{pytest},
    GitHub Actions).
  \item Use a larger or real street map to retest the 80\% target.
  \item Send real SOS messages, run a pilot with 10--20 peers, load test
    and deploy.
\end{enumerate}

\section{Final Remarks}

The prototype shows that making safety a routing input is feasible with
modest tools. The main lesson is to get one complete loop working early
and let honest measurements, including the missed target, decide what to
improve next.

\printbibliography[heading=bibintoc]

\end{document}